\suggestions


\chapter*{RECOMENDACIONES}
\addcontentsline{toc}{chapter}{RECOMENDACIONES}

A partir de los resultados obtenidos y las experiencias derivadas del desarrollo de esta investigación, se proponen las siguientes recomendaciones para el perfeccionamiento y evolución de Assets Studio:

\begin{enumerate}
    \item Evaluar la integración de modelos de Inteligencia Artificial generativa de ejecución local (\textit{Edge AI}) para permitir la creación y el etiquetado de recursos básicos incluso en situaciones de ausencia total de conectividad, reduciendo la dependencia de APIs externas.

    \item Expandir el sistema de \textit{Scraping} Ético hacia un mayor número de repositorios de recursos abiertos, implementando algoritmos de filtrado más granulares que permitan clasificar los activos no solo por tipo, sino por compatibilidad específica con motores de renderizado y \textit{software} de diseño profesional.

    \item Implementar un módulo de analítica avanzada para administradores de equipos que permita visualizar patrones de uso de los activos, ayudando a identificar cuáles son los recursos más demandados y facilitando la toma de decisiones sobre la adquisición o creación de nuevo material creativo.

    \item Desarrollar extensiones o \textit{plugins} nativos para herramientas líderes del sector (como Adobe Creative Cloud, Blender o VS Code), de manera que los usuarios puedan acceder al catálogo centralizado de Assets Studio directamente desde su entorno de trabajo habitual sin necesidad de cambiar de aplicación.

    \item Fomentar la creación de una comunidad de colaboradores que contribuya al entrenamiento continuo del modelo de búsqueda semántica, mediante la validación manual de etiquetas sugeridas por la IA, asegurando así que el sistema evolucione con el lenguaje técnico y las tendencias del sector creativo.
\end{enumerate}